\documentclass[tikz,border=6pt]{standalone}
% Shared visual language for the thesis architecture figures.
% Compile each standalone source with XeLaTeX from this directory.
\usepackage[UTF8,scheme=plain,fontset=none]{ctex}
\usepackage{amsmath,amssymb}
\usepackage{fontspec}
\setmainfont{TeX Gyre Termes}
\setsansfont{TeX Gyre Heros}
\setmonofont{DejaVu Sans Mono}
\setCJKmainfont{Noto Serif SC}
\setCJKsansfont[BoldFont=SimHei]{Noto Sans SC}
\usetikzlibrary{arrows.meta,positioning,calc,fit,backgrounds,shapes.geometric,matrix,decorations.pathreplacing}

\definecolor{Ink}{HTML}{203040}
\definecolor{Muted}{HTML}{66788A}
\definecolor{LineSoft}{HTML}{AAB7C4}
\definecolor{DataBlue}{HTML}{2563A6}
\definecolor{DataFill}{HTML}{EAF3FB}
\definecolor{ClockPurple}{HTML}{7450A5}
\definecolor{ClockFill}{HTML}{F1ECF8}
\definecolor{ControlTeal}{HTML}{16827B}
\definecolor{ControlFill}{HTML}{E9F6F4}
\definecolor{DspPurple}{HTML}{6A4C93}
\definecolor{DspFill}{HTML}{F0EAF7}
\definecolor{RamGreen}{HTML}{367C59}
\definecolor{RamFill}{HTML}{ECF6F0}
\definecolor{FabricOrange}{HTML}{C56A1A}
\definecolor{FabricFill}{HTML}{FFF1E5}
\definecolor{EvidenceGreen}{HTML}{2F7D4B}
\definecolor{EvidenceFill}{HTML}{EAF6EE}
\definecolor{Amber}{HTML}{B07A16}
\definecolor{AmberFill}{HTML}{FFF6DD}
\definecolor{PendingGray}{HTML}{7C8792}
\definecolor{PendingFill}{HTML}{F1F3F5}
\definecolor{Danger}{HTML}{A04444}
\definecolor{DangerFill}{HTML}{FFF0F0}

\tikzset{
  >={Stealth[length=2.2mm,width=1.45mm]},
  every node/.style={font=\sffamily\footnotesize,text=Ink},
  data/.style={-{Stealth[length=2.2mm,width=1.45mm]},draw=DataBlue,line width=0.92pt},
  data bi/.style={{Stealth[length=1.9mm,width=1.3mm]}-{Stealth[length=1.9mm,width=1.3mm]},draw=DataBlue,line width=0.82pt},
  control/.style={-{Stealth[length=2.0mm,width=1.35mm]},draw=ControlTeal,dashed,line width=0.78pt},
  clock/.style={-{Stealth[length=2.0mm,width=1.35mm]},draw=ClockPurple,densely dotted,line width=0.9pt},
  address/.style={-{Stealth[length=2.0mm,width=1.35mm]},draw=FabricOrange,dash dot,line width=0.76pt},
  feedback/.style={-{Stealth[length=1.8mm,width=1.25mm]},draw=Muted,line width=0.66pt},
  block/.style={draw=Ink,fill=white,rounded corners=1.5pt,line width=0.68pt,minimum height=10mm,minimum width=23mm,align=center,inner sep=3pt},
  data block/.style={block,draw=DataBlue,fill=DataFill},
  control block/.style={block,draw=ControlTeal,fill=ControlFill},
  clock block/.style={block,draw=ClockPurple,fill=ClockFill},
  dsp/.style={block,draw=DspPurple,fill=DspFill,line width=0.9pt},
  ram/.style={block,draw=RamGreen,fill=RamFill,line width=0.9pt},
  fabric/.style={block,draw=FabricOrange,fill=FabricFill,line width=0.9pt},
  result/.style={block,draw=EvidenceGreen,fill=EvidenceFill,line width=0.9pt},
  pending/.style={block,draw=PendingGray,fill=PendingFill,line width=0.8pt},
  warning/.style={block,draw=Amber,fill=AmberFill,line width=0.9pt},
  boundary/.style={draw=Muted!75,rounded corners=2.2mm,dashed,line width=0.55pt,inner sep=3mm},
  title/.style={font=\sffamily\bfseries\large,text=Ink,align=center},
  subtitle/.style={font=\sffamily\scriptsize,text=Muted,align=center},
  group title/.style={font=\sffamily\bfseries\footnotesize,text=Ink,anchor=west},
  signal/.style={font=\sffamily\scriptsize,text=Ink,fill=white,inner sep=1.2pt,align=center},
  note/.style={font=\sffamily\scriptsize,text=Muted,align=center},
  tiny note/.style={font=\sffamily\tiny,text=Muted,align=center},
  badge/.style={draw=EvidenceGreen,fill=EvidenceFill,rounded corners=1.8mm,line width=0.65pt,font=\sffamily\bfseries\scriptsize,text=EvidenceGreen,inner xsep=4pt,inner ysep=2pt,align=center},
  pending badge/.style={draw=PendingGray,fill=PendingFill,rounded corners=1.8mm,line width=0.65pt,font=\sffamily\bfseries\scriptsize,text=PendingGray,inner xsep=4pt,inner ysep=2pt,align=center},
  innovation/.style={circle,draw=Amber,fill=AmberFill,line width=0.8pt,minimum size=6.5mm,font=\sffamily\bfseries\scriptsize,text=Amber,inner sep=0pt},
  sum/.style={circle,draw=Ink,fill=white,line width=0.75pt,minimum size=6mm,inner sep=0pt},
  mux/.style={trapezium,trapezium left angle=72,trapezium right angle=108,draw=Ink,fill=white,line width=0.68pt,minimum height=11mm,minimum width=12mm,align=center,inner sep=2pt}
}

\newcommand{\bitrate}[2]{{\scriptsize #1 bit，#2}}
\newcommand{\passmark}{\textcolor{EvidenceGreen}{\bfseries 满足判据}}
\newcommand{\pendingmark}{\textcolor{PendingGray}{\bfseries 待验证}}


\begin{document}
\begin{tikzpicture}[x=1cm,y=1cm]

\node[title] at (12.5,12.15) {当前2-DSP目标配置的资源角色、器件占用与Pareto位置};
\node[subtitle] at (12.5,11.70)
  {Vivado 2025.2，XC7A35T-FGG484-2，完整板级布局布线后，24/20/20 bit，CIC mode 0};

% Exact board total card.
\node[result,minimum width=220mm,minimum height=15mm] (totals) at (12.50,10.65)
  {\bfseries 268 LUT \quad 417 FF \quad 2 DSP48E1 \quad 4 RAMB18E1（2 BRAM Tile）\quad 2 MMCM\qquad
   WNS/WHS=+44.389/+0.078 ns，TNS/THS=0/0，DRC Error=0，配置位流生成完成};
\node[pending badge] at (23.45,11.43) {工具验证完成\\物理板待验证};

% Structural resource roles.
\node[dsp,minimum width=34mm,minimum height=17mm] (dsp0) at (2.35,8.05)
  {DSP48E1\#0\\[-1pt]{\scriptsize Stage1串行MAC}};
\node[dsp,minimum width=37mm,minimum height=17mm] (dsp1) at (6.55,8.05)
  {DSP48E1\#1\\[-1pt]{\scriptsize Stage2/3时分MAC}};
\node[fabric,minimum width=42mm,minimum height=17mm] (cic) at (11.25,8.05)
  {Fabric/CARRY + FF\\[-1pt]{\scriptsize CIC差分、Hold、26/29 bit积分}\\[-1pt]{\scriptsize CIC占用DSP=0}};
\node[ram,minimum width=30mm,minimum height=17mm] (rh1) at (2.15,5.95)
  {RAMB18E1\#1\\[-1pt]{\scriptsize Stage1历史}};
\node[ram,minimum width=34mm,minimum height=17mm] (rh23) at (5.85,5.95)
  {RAMB18E1\#2\\[-1pt]{\scriptsize Stage2/3历史}};
\node[ram,minimum width=33mm,minimum height=17mm] (rc) at (9.75,5.95)
  {RAMB18E1\#3\\[-1pt]{\scriptsize 统一FIR系数}};
\node[pending,minimum width=31mm,minimum height=17mm] (rrom) at (13.45,5.95)
  {RAMB18E1\#4\\[-1pt]{\scriptsize 板级测试音ROM}};
\node[clock block,minimum width=35mm,minimum height=17mm] (rclk) at (6.90,3.95)
  {2$\times$MMCME2\_ADV\\[-1pt]{\scriptsize 44.1/48 kHz双时钟族}};
\node[pending,minimum width=47mm,minimum height=17mm] (rest) at (12.10,3.95)
  {其余板级逻辑\\[-1pt]{\scriptsize 键盘、CDC、CE、输出MUX与DAC接口}};
\draw[data] (rh1.north) -- (dsp0.south);
\draw[data] (rh23.north) -- (dsp1.south);
\draw[data] (rc.north) -- ++(0,0.40) -| (dsp0.south east);
\draw[data] (rc.north) -- ++(0,0.40) -| (dsp1.south east);

\begin{scope}[on background layer]
  \node[boundary,fill=DataFill!28,fit=(dsp0)(dsp1)(cic)(rh1)(rh23)(rc)] (coregrp) {};
\end{scope}
\node[group title] at ($(coregrp.north west)+(0.20,-0.25)$) {滤波器核心的明确资源角色：2 DSP + 3 RAMB18};
\node[font=\sffamily\bfseries\footnotesize] at (10.20,4.95) {核心外板级资源};

% Utilization bars: two explicit scales to preserve readability.
\node[font=\sffamily\bfseries\footnotesize,anchor=west] at (16.20,9.25) {器件占用率（完整板级）};
\node[tiny note,anchor=west] at (16.20,8.88) {逻辑/DSP/BRAM/IO横条量程0--10\%};
\newcommand{\utilbar}[5]{%
  \node[tiny note,anchor=east,text=Ink] at (19.10,#1) {#2};%
  \draw[LineSoft,rounded corners=1pt,line width=0.45pt] (19.28,#1-0.12) rectangle ++(4.80,0.24);%
  \fill[#5,rounded corners=1pt] (19.28,#1-0.12) rectangle ++(#4,0.24);%
  \node[tiny note,anchor=west,text=Ink] at (24.22,#1) {#3};%
}
\utilbar{8.38}{LUT 268/20800}{1.29\%}{0.6192}{DataBlue}
\utilbar{7.78}{FF 417/41600}{1.00\%}{0.4800}{ControlTeal}
\utilbar{7.18}{DSP 2/90}{2.22\%}{1.0656}{DspPurple}
\utilbar{6.58}{BRAM Tile 2/50}{4.00\%}{1.9200}{RamGreen}
\utilbar{5.98}{IOB 17/250}{6.80\%}{3.2640}{FabricOrange}
\node[tiny note,anchor=west] at (16.20,5.35) {MMCM横条量程0--50\%};
\node[tiny note,anchor=east,text=Ink] at (19.10,4.88) {MMCM 2/5};
\draw[LineSoft,rounded corners=1pt,line width=0.45pt] (19.28,4.76) rectangle ++(4.80,0.24);
\fill[ClockPurple,rounded corners=1pt] (19.28,4.76) rectangle ++(3.84,0.24);
\node[tiny note,anchor=west,text=Ink] at (24.22,4.88) {40.00\%};
\node[note,anchor=west,align=left] at (16.20,4.15)
  {MMCM占用率较高源于器件仅含5个MMCM，\\其绝对用量为2，并非数据通路的主要算术成本。};

% Pareto staircase.
\node[font=\sffamily\bfseries\footnotesize] at (12.50,2.95) {DSP--LUT Pareto端点（RAMB18E1均为4个）};
\node[result,minimum width=48mm,minimum height=18mm] (p218) at (4.10,1.70)
  {LUT优先端点\\[-1pt]{\bfseries 218 LUT / 365 FF / 4 DSP}\\[-1pt]{\scriptsize 已完成物理板功能验证}};
\node[result,minimum width=48mm,minimum height=18mm] (p239) at (12.50,1.70)
  {资源折中端点\\[-1pt]{\bfseries 239 LUT / 388 FF / 3 DSP}\\[-1pt]{\scriptsize 已完成物理板功能验证}};
\node[pending,minimum width=48mm,minimum height=18mm] (p268) at (20.90,1.70)
  {DSP受限目标配置\\[-1pt]{\bfseries 268 LUT / 417 FF / 2 DSP}\\[-1pt]{\scriptsize 工具验证完成，物理板待验证}};
\draw[data] (p218) -- node[signal,above] {$+21$ LUT，$+23$ FF，$-1$ DSP} (p239);
\draw[data] (p239) -- node[signal,above] {$+29$ LUT，$+29$ FF，$-1$ DSP} (p268);

\node[note,anchor=west] at (0.35,0.18)
  {注：模块框仅表示可由RTL/原语实例明确确定的资源角色；布局布线后扁平网表中的LUT与FF不按模块人为分摊。};
\node[note,anchor=east] at (24.80,-0.18)
  {268版相对239版以29 LUT和29 FF替代1个DSP，差异与末级积分器由DSP改映射为Fabric一致。};

\end{tikzpicture}
\end{document}


